\subsection{\ifzh 低海拔网格的占域持续概率随绿度而异\else Low-elevation occupancy persistence varies with greenness\fi}
\ifzh
将校正后的转移概率投射到实际地点环境，得到比单独报告系数更明确的持续性结构（表~\ref{tab:envpersistence}；图~\ref{fig:envpersistence}）。低海拔低、中、高绿度地点的平均占域持续概率分别为 24.2\%、49.7\% 和 71.3\%，对应 320、430 和 158 个地点；高海拔高绿度的 39 个地点为 96.8\%。这说明持续性梯度不仅体现在山地与低地之间，低地内部也存在明显的绿度相关分化。与报告分层结果对照，高海拔高绿度环境同时具有较高报告水平和模型估计的较高占域持续概率，支持在这些环境中设置固定调查网格，检验后续时期的占域状态。在低海拔高绿度环境中，报告概率的点估计尚未回到初期水平，但地点占域持续概率相对较高。报告回升程度与已占据地点的保留描述不同过程，低地整体报告偏低不能直接推导为所有低地地点持续性较弱。这些分层概率由同一校正模型推导，提供的是环境关联的具体表达。
\else
Projecting corrected transition probabilities onto observed site environments reveals a persistence structure beyond the coefficient signs (Table~\ref{tab:envpersistence}; Figure~\ref{fig:envpersistence}). Mean occupancy persistence probability was 24.2\%, 49.7\% and 71.3\% across low-elevation low-, middle- and high-greenness cells, containing 320, 430 and 158 sites. The 39 high-elevation, high-greenness sites averaged 96.8\%. Persistence therefore varies both between uplands and lowlands and along greenness within lowlands. Compared with the reporting strata, high-elevation, high-greenness environments combine high reporting with high modeled persistence, identifying environments in which fixed survey grids can test occupancy in later periods. In greener lowlands, the reporting point estimate remains below its initial baseline despite relatively high occupancy persistence probability. Rebound and retention of occupied sites describe different processes, so low aggregate reporting does not imply uniformly weak persistence across lowlands. These probabilities express environmental associations derived from the same corrected model.
\fi
\inctab{environment_persistence}
\ifzh
固定海拔在 100~m 后，NDVI 从 0.30 增至 0.76 时，占域持续概率预测由 29.8\% 升至 72.0\%，差值为 42.2 个百分点（95\% 区间 9.9--64.9）；定殖预测仅由 5.3\% 变为 6.0\%，差值 0.7 个百分点，区间为 $[-2.9,5.3]$。该对照保持海拔不变，支持绿度主要对应已占据地点保留的解释。对保护管理而言，维持已有占域地点的植被条件，与促进未占据地点进入占域状态，是需要分别追踪的目标；较绿环境中较强的保留关联不能直接换算为植被增加后的定殖收益。
\else
Holding elevation at 100~m, increasing NDVI from 0.30 to 0.76 raised predicted occupancy persistence probability from 29.8\% to 72.0\%, a 42.2-point contrast (95\% interval 9.9--64.9). Colonization changed from 5.3\% to 6.0\%, a 0.7-point contrast with interval $[-2.9,5.3]$. Holding elevation fixed supports an interpretation in which greenness chiefly characterizes retention of occupied sites. Maintaining vegetation conditions at occupied sites and promoting transitions from nonoccupancy to occupancy are therefore distinct monitoring targets; the retention association does not directly quantify a colonization benefit from increasing vegetation.
\fi
\begin{figure}[htbp]\centering
\incfig{fig27_environment_persistence}
\caption{\ifzh
校正模型与等访问次数的报告历史。（a）占域持续概率；（b）定殖；（c）前期有报告后再次报告的比例及计数。前两图标注地点数，阴影表示不足十个地点的组合。过程概率对应模型主要时期之间的转移。
\else
Corrected-model processes and equal-visit reporting histories. Panels show (a) occupancy persistence probability, (b) colonization, and (c) repeated reports conditional on a previous report, with counts. Site counts appear in the first two panels; hatching denotes fewer than ten sites. Process probabilities concern transitions between model primary periods.
\fi}\label{fig:envpersistence}
\end{figure}
\ifzh
等访问次数的记录也显示出环境差异。低海拔低绿度有 1/9 个前期报告地点--时期对再次报告，高绿度为 15/28；高海拔高绿度为 22/32。较绿环境中较高的再次记录比例，与模型估计的较高占域持续概率方向一致；再次记录仍同时受到物种存在与探测机会影响。例如，高海拔高绿度中，前期未报告的十个完整访问对有七个在下一时期报告，而模型定殖概率约为 6\%。较高的后期记录比例与较低的定殖估计并列，说明未报告地点中包含模型判断仍可能已占据的地点。重复调查形成的探测历史有助于区分已占据网格中的再次探测和由未占据到占据的状态转移，从而解释飓风后的报告回升。
\else
Equal-visit histories also differed across environments. Repeated reports occurred in 1/9 previously reporting low-elevation, low-greenness period pairs, 15/28 low-elevation, high-greenness pairs, and 22/32 high-elevation, high-greenness pairs. Higher repeat-report proportions in greener environments are directionally consistent with higher modeled occupancy persistence; repeat reports depend on both presence and detection. Among ten fully sampled high-elevation, high-greenness pairs with no earlier report, seven had a later report, despite modeled colonization near 6\%. Frequent later reports among previously nonreporting pairs alongside low modeled colonization illustrates that nonreporting sites can still be occupied under the detection model. Repeated surveys are therefore valuable for distinguishing renewed detection at occupied grids from transitions out of nonoccupancy.
\fi
\FloatBarrier
